% GENERATED FILE. Edit canonical lessons, then run npm run generate:books.
% canonical-source-hash: fnv1a64:2cdb48ed0738979c
% canonical-lessons: TE-C420-empika, TE-C420-pratyamnayam, TE-C420-pariganincu, TE-C420-aicchikam, TE-C420-vaduluko

\chapter{Making Choices}
\label{ch:te-making-choices}
\hlchaptermodality{420}

\begin{chapteropening}
\textbf{By the end of this chapter:} \emph{I can talk about choices, alternatives and giving things up.}

Complete the last lesson of chapter 420: I can talk about choices, alternatives and giving things up.
\end{chapteropening}

\section[\texorpdfstring{\te{ఎంపిక} (empika)}{ఎంపిక (empika)}]{\te{ఎంపిక} (empika) — a choice; a selection}

\label{lesson:TE-C420-empika}



\begin{quote}
Before the new one: say the Telugu for in any case, then the Telugu for to occur.
\end{quote}

\subsection*{You'll want to know: \te{ఎంపిక}}
\textbf{\te{ఎంపిక}} — \emph{empika} — \textquotedblleft{}a choice; a selection\textquotedblright{}.

\begin{quote}
\te{ఆ} \te{ఉద్యోగానికి} \te{ఆమెను} \te{తీసుకోవడం} \te{మంచి} \te{ఎంపిక}.
\end{quote}

\begin{quote}
\emph{ā udyōgāniki āmenu tīsukōvaḍaṁ mañci empika.}
\end{quote}

\begin{quote}
Taking her on for that job was a good choice.
\end{quote}

\te{ఎంపిక} \te{చేయు}, to select; \te{ఎంపికైన}, selected.

\begin{cousinweb}
\textbf{\te{ఎంపిక}} (\emph{empika}) is a Telugu noun whose origin is not settled; DEDR does not list it.
\end{cousinweb}

\begin{grammarlens}[title={What you've built}]
One more word for saying what happens if.
\end{grammarlens}

\subsection*{Your turn}
\begin{itemize}
\raggedright
  \item \emph{Say it:} \emph{empika}
  \item \emph{Say it:} \emph{empika}, once more
  \item \emph{Say it:} \emph{ā udyōgāniki āmenu tīsukōvaḍaṁ mañci empika.} — Taking her on for that job was a good choice.
  \item \emph{Recall:} say the Telugu for in any case, then the Telugu for to occur, then say \emph{empika} again
\end{itemize}

\begin{tcolorbox}[breakable,colback=teal!4,colframe=teal!35!black,title={Before you move on}]
What does \te{ఎంపిక} mean? (\textquotedblleft{}A choice; a selection\textquotedblright{}.) Say it once more.
\end{tcolorbox}
\section[\texorpdfstring{\te{ప్రత్యామ్నాయం} (pratyāmnāyaṁ)}{ప్రత్యామ్నాయం (pratyāmnāyaṁ)}]{\te{ప్రత్యామ్నాయం} (pratyāmnāyaṁ) — an alternative}

\label{lesson:TE-C420-pratyamnayam}



\begin{quote}
Before the new one: say the Telugu for to occur, then the Telugu for a choice.
\end{quote}

\subsection*{You'll want to know: \te{ప్రత్యామ్నాయం}}
\textbf{\te{ప్రత్యామ్నాయం}} — \emph{pratyāmnāyaṁ} — \textquotedblleft{}an alternative\textquotedblright{}.

\begin{quote}
\te{మాకు} \te{వేరే} \te{ప్రత్యామ్నాయం} \te{లేదు}.
\end{quote}

\begin{quote}
\emph{māku vērē pratyāmnāyaṁ lēdu.}
\end{quote}

\begin{quote}
We have no other alternative.
\end{quote}

As an adjective: \te{ప్రత్యామ్నాయ} \te{ఏర్పాటు}, an alternative arrangement.

\begin{cousinweb}
\textbf{\te{ప్రత్యామ్నాయం}} (\emph{pratyāmnāyaṁ}) is borrowed from Sanskrit \emph{pratyāmnāya}, a term of logic for \textquotedblleft{}the proposition restated\textquotedblright{} (Monier-Williams); the Telugu sense \textquotedblleft{}an alternative\textquotedblright{} is modern.
\end{cousinweb}

\begin{grammarlens}[title={What you've built}]
One more word for saying what happens if.
\end{grammarlens}

\subsection*{Your turn}
\begin{itemize}
\raggedright
  \item \emph{Say it:} \emph{pratyāmnāyaṁ}
  \item \emph{Say it:} \emph{pratyāmnāyaṁ}, once more
  \item \emph{Say it:} \emph{māku vērē pratyāmnāyaṁ lēdu.} — We have no other alternative.
  \item \emph{Recall:} say the Telugu for to occur, then the Telugu for a choice, then say \emph{pratyāmnāyaṁ} again
\end{itemize}

\begin{tcolorbox}[breakable,colback=teal!4,colframe=teal!35!black,title={Before you move on}]
What does \te{ప్రత్యామ్నాయం} mean? (\textquotedblleft{}An alternative\textquotedblright{}.) Say it once more.
\end{tcolorbox}
\section[\texorpdfstring{\te{పరిగణించు} (parigaṇiñcu)}{పరిగణించు (parigaṇiñcu)}]{\te{పరిగణించు} (parigaṇiñcu) — to take into account}

\label{lesson:TE-C420-pariganincu}



\begin{quote}
Before the new one: say the Telugu for a choice, then the Telugu for an alternative.
\end{quote}

\subsection*{You'll want to know: \te{పరిగణించు}}
\textbf{\te{పరిగణించు}} — \emph{parigaṇiñcu} — \textquotedblleft{}to take into account\textquotedblright{}.

\begin{quote}
\te{అతని} \te{వయసును} \te{పరిగణించి} \te{అతన్ని} \te{క్షమించారు}.
\end{quote}

\begin{quote}
\emph{atani vayasunu parigaṇiñci atanni kṣamiñcāru.}
\end{quote}

\begin{quote}
Taking his age into account, they forgave him.
\end{quote}

Past (I …): \te{పరిగణించాను}, \emph{parigaṇiñcānu}.

\te{పరిగణనలోకి} \te{తీసుకో}, to take into consideration.

\begin{cousinweb}
\textbf{\te{పరిగణించు}} (\emph{parigaṇiñcu}) is built on Sanskrit \emph{parigaṇana} (Monier-Williams), \emph{pari-}, \textquotedblleft{}all round\textquotedblright{}, with \emph{gaṇana}, \textquotedblleft{}counting\textquotedblright{} (CDIAL 3992), with the Telugu verb ending -\te{ఇంచు}.
\end{cousinweb}

\begin{grammarlens}[title={What you've built}]
One more word for saying what happens if.
\end{grammarlens}

\subsection*{Your turn}
\begin{itemize}
\raggedright
  \item \emph{Say it:} \emph{parigaṇiñcu}
  \item \emph{Say it:} \emph{parigaṇiñcu}, once more
  \item \emph{Say it:} \emph{atani vayasunu parigaṇiñci atanni kṣamiñcāru.} — Taking his age into account, they forgave him.
  \item \emph{Recall:} say the Telugu for a choice, then the Telugu for an alternative, then say \emph{parigaṇiñcu} again
\end{itemize}

\begin{tcolorbox}[breakable,colback=teal!4,colframe=teal!35!black,title={Before you move on}]
What does \te{పరిగణించు} mean? (\textquotedblleft{}To take into account\textquotedblright{}.) Say it once more.
\end{tcolorbox}
\section[\texorpdfstring{\te{ఐచ్ఛికం} (aicchikaṁ)}{ఐచ్ఛికం (aicchikaṁ)}]{\te{ఐచ్ఛికం} (aicchikaṁ) — optional}

\label{lesson:TE-C420-aicchikam}



\begin{quote}
Before the new one: say the Telugu for an alternative, then the Telugu for to take into account.
\end{quote}

\subsection*{You'll want to know: \te{ఐచ్ఛికం}}
\textbf{\te{ఐచ్ఛికం}} — \emph{aicchikaṁ} — \textquotedblleft{}optional\textquotedblright{}.

\begin{quote}
\te{ఆ} \te{పాఠం} \te{ఐచ్ఛికం}, \te{తప్పనిసరి} \te{కాదు}.
\end{quote}

\begin{quote}
\emph{ā pāṭhaṁ aicchikaṁ, tappanisari kādu.}
\end{quote}

\begin{quote}
That lesson is optional, not compulsory.
\end{quote}

The opposite of \te{తప్పనిసరి}, compulsory, which you know.

\begin{cousinweb}
\textbf{\te{ఐచ్ఛికం}} (\emph{aicchikaṁ}) is borrowed from Sanskrit \emph{aicchika}, \textquotedblleft{}optional, arbitrary, at will\textquotedblright{} (Monier-Williams), from \emph{icchā}, \textquotedblleft{}wish\textquotedblright{}.
\end{cousinweb}

\begin{grammarlens}[title={What you've built}]
One more word for saying what happens if.
\end{grammarlens}

\subsection*{Your turn}
\begin{itemize}
\raggedright
  \item \emph{Say it:} \emph{aicchikaṁ}
  \item \emph{Say it:} \emph{aicchikaṁ}, once more
  \item \emph{Say it:} \emph{ā pāṭhaṁ aicchikaṁ, tappanisari kādu.} — That lesson is optional, not compulsory.
  \item \emph{Recall:} say the Telugu for an alternative, then the Telugu for to take into account, then say \emph{aicchikaṁ} again
\end{itemize}

\begin{tcolorbox}[breakable,colback=teal!4,colframe=teal!35!black,title={Before you move on}]
What does \te{ఐచ్ఛికం} mean? (\textquotedblleft{}Optional\textquotedblright{}.) Say it once more.
\end{tcolorbox}
\section[\texorpdfstring{\te{వదులుకో} (vadulukō)}{వదులుకో (vadulukō)}]{\te{వదులుకో} (vadulukō) — to let go of; to give up}

\label{lesson:TE-C420-vaduluko}



\begin{quote}
Before the new one: say the Telugu for to take into account, then the Telugu for optional.
\end{quote}

\subsection*{You'll want to know: \te{వదులుకో}}
\textbf{\te{వదులుకో}} — \emph{vadulukō} — \textquotedblleft{}to let go of; to give up\textquotedblright{}.

\begin{quote}
\te{అతను} \te{ఆ} \te{ఉద్యోగం} \te{కోసం} \te{తన} \te{చదువు} \te{వదులుకున్నాడు}.
\end{quote}

\begin{quote}
\emph{atanu ā udyōgaṁ kōsaṁ tana caduvu vadulukunnāḍu.}
\end{quote}

\begin{quote}
He gave up his studies for that job.
\end{quote}

Past (I …): \te{వదులుకున్నాను}, \emph{vadulukunnānu}.

\te{వదులు}, to let loose, + -\te{కో}: to let a thing go from oneself.

\begin{cousinweb}
\textbf{\te{వదులుకో}} (\emph{vadulukō}) is native Dravidian: DEDR 5243 lists Telugu \emph{vadalu}, \textquotedblleft{}to leave, to abandon, to let loose\textquotedblright{}; -\te{కో}, from \emph{konu}, \textquotedblleft{}to take\textquotedblright{} (DEDR 2151), turns it back on the doer.
\end{cousinweb}

\begin{grammarlens}[title={What you've built}]
One more word for saying what happens if.
\end{grammarlens}

\subsection*{Your turn}
\begin{itemize}
\raggedright
  \item \emph{Say it:} \emph{vadulukō}
  \item \emph{Say it:} \emph{vadulukō}, once more
  \item \emph{Say it:} \emph{atanu ā udyōgaṁ kōsaṁ tana caduvu vadulukunnāḍu.} — He gave up his studies for that job.
  \item \emph{Recall:} say the Telugu for to take into account, then the Telugu for optional, then say \emph{vadulukō} again
\end{itemize}

\begin{tcolorbox}[breakable,colback=teal!4,colframe=teal!35!black,title={Before you move on}]
What does \te{వదులుకో} mean? (\textquotedblleft{}To let go of; to give up\textquotedblright{}.) Say it once more.
\end{tcolorbox}
